% ==============================================================
\section{Experiments}\label{sec:experiments}
% ==============================================================

We verify the scalability and effectiveness of the proposed method.
All experiments were implemented in Rust (release build)
and executed on Apple M3 (8 cores, 24\,GB RAM).

% ==============================================================
\subsection{EX-1: Method Comparison on Diverse Dense Region Shapes (Synthetic Data)}
\label{sec:ex1}\label{sec:ex1_1}
% ==============================================================

\paragraph{Objective.}
Real-world dense regions may take various shapes.
This experiment uses synthetic data with four embedding shapes
(rectangular, drift, intermittent, and expanding)
to verify how well existing methods recover the exact solution
defined by the proposed method.

\paragraph{Synthetic data.}
A $T{=}100 \times S{=}100$ grid with 20 embedded GT patterns is used.
Four embedding shapes are evaluated:
rectangular (Shape~A), drift (Shape~B), intermittent (Shape~C),
and expanding (Shape~D),
with 3 seeds per shape ($4 \times 3 = 12$ datasets in total).
Generation details are given in Appendix~\ref{sec:ex1_data_details}.

\paragraph{Baseline methods.}
Each baseline addresses a fundamentally different formulation,
so all tunable parameters were grid-searched to maximize each baseline's score for a fair comparison.
The following five methods are used.
Colocation / STCP-Miner are used in EX-1-1 (pattern detection);
LISA / MZoom / ST-DBSCAN are used in EX-1-2 (region recovery).

\begin{itemize}
\item \textbf{Colocation}~\cite{shekhar2001colocation}: Spatial colocation mining. PI threshold $\tau$ grid-searched.
\item \textbf{STCP-Miner}~\cite{qian2009stcolocation}: Spatio-temporal window extension of Colocation. $\tau$ and window width grid-searched.
\item \textbf{MZoom}~\cite{shin2018mzoom}: Greedy dense block detection on the occurrence tensor of GT pattern $P$. Top $k$ grid-searched.
\item \textbf{LISA}~\cite{anselin1995lisa}: Occurrence counts of GT pattern $P$ aggregated onto a spatial grid; High-High cells identified by Local Moran's $I$.
\item \textbf{ST-DBSCAN}~\cite{birant2007stdbscan}: $\varepsilon_1$, $\varepsilon_2$, min\_pts grid-searched on occurrence cells of GT pattern $P$.
\end{itemize}

\paragraph{Evaluation metrics.}
The dense pattern set $P^*$ and dense regions $R^*(P)$
obtained under the sliding-window density definition (Definitions~3--5)
for the explicitly embedded patterns serve as the ground truth (GT).
Since the proposed method outputs the exact solution of this definition,
it achieves F1 $= 1.00$ and Jaccard $= 1.00$ by construction.
For EX-1-1,
F1 $= 2|\hat{P} \cap P^*| / (|\hat{P}| + |P^*|)$
measures how well each baseline identifies the GT pattern set $P^*$.
For EX-1-2, mean Jaccard
$= \frac{1}{|P^*|}\sum_{P \in P^*} |R^*(P) \cap \hat{R}(P)| / |R^*(P) \cup \hat{R}(P)|$
measures how closely each baseline recovers the GT dense region $R^*(P)$.

Note that each baseline solves a fundamentally different problem:
colocation methods detect participation-index co-occurrence without defining dense regions;
dense subtensor methods output Cartesian-product blocks, not sliding-window regions;
clustering methods label point neighborhoods rather than windowed occurrence counts.
EX-1 is therefore not a head-to-head competition but a characterization of
how closely each method approximates the exact sliding-window solution.

\paragraph{Proposed method parameters.}
Common across all shapes: $W = (20, 5)$, $\theta = 30$.

\paragraph{Results.}
Fig.~\ref{fig:ex1} shows scores averaged over 3 seeds per shape.

\begin{figure}[t]
\centering
\includegraphics[width=\columnwidth]{fig/fig_ex1_heatmap.pdf}
\caption{EX-1 results (3-seed average, $p_{\mathrm{embed}}{=}0.25$).
  (a)~EX-1-1 F1; (b)~EX-1-2 mean Jaccard.
  GT is the exact solution under the sliding-window definition (Definitions~3--5).
  $\dagger$The proposed method outputs the exact solution by construction,
  hence F1 $= 1.00$ and Jaccard $= 1.00$.
  Baseline parameters selected by grid search.}
\label{fig:ex1}
\end{figure}

\paragraph{Analysis.}
\emph{EX-1-1 (Pattern detection F1).}
Colocation performs worst on Shape~D (expanding):
during the expansion phase, the global PI computation is diluted
over the entire observation period, suppressing the participation index.
STCP-Miner partially avoids this dilution via temporal windowing,
but the approximation gap remains.

\emph{EX-1-2 (Region recovery, mean Jaccard).}
All baselines remain at low scores across all shapes.

\textbf{LISA} collapses occurrences over all time steps, irreversibly discarding temporal structure.

\textbf{MZoom} achieves the lowest scores.
Against an $S{=}100$-cell grid where each dense band spans only 5 cells,
any Cartesian-product block wide enough to cover the full band
incurs a severe density drop from the Bernoulli background.
Greedy search therefore converges to single-column (width~1) blocks,
missing most of the true dense region as false negatives.

\textbf{ST-DBSCAN} labels only occurrence cells as cluster members,
systematically missing non-occurrence cells inside dense windows,
which constitute the majority of the dense region.

In short, these structural mismatches---not parameter tuning---explain
why no baseline approximates the sliding-window definition well
across all four shapes.

% ==============================================================
\subsection{EX-2: Scalability (Synthetic Data)}
\label{sec:ex2}
% ==============================================================

\paragraph{Objective.}
We verify how the computation time of the proposed method
scales with increasing number of transactions $N$,
and whether the ratio of the number of temporal grid points $T$
to the number of spatial grid points $S$ affects performance
at the same $N$ (for $d=2$, $N = T \times S$).

\paragraph{Synthetic data generation.}
On a $d = 2$ lattice
$\{0,\ldots,T{-}1\} \times \{0,\ldots,S{-}1\}$,
each item is independently sampled from
Bernoulli$(p_{\mathrm{item}}{=}0.10)$
for all cells to form one basket.
Empty cells are assigned item~1 by default,
ensuring exactly one transaction per cell,
so that $N = T \times S$
(generation details in Appendix~\ref{sec:ex2_data_details}).

\paragraph{Proposed method parameters.}
$|\mathcal{I}| = 20$, $W = (24, 4)$, and $\theta = 15$ were fixed.

\paragraph{Rationale for experimental axes.}
Since linear scaling with $|\mathcal{I}|$ and $p_{\mathrm{item}}$ was confirmed
in the 1D version~\cite{aprioriwindow},
and the effect of $d$ is absorbed into $N = T \times S^{d-1}$,
we fix $d = 2$ and vary $T$ and $S$.

\paragraph{Experimental design.}
$S \in \{10,\; 100,\; 1{,}000,\; 10{,}000\}$ and
$T \in \{100,\; 1{,}000,\; 10{,}000,\; 100{,}000\}$
yield $4 \times 4 = 16$ configurations.
$N = T \times S$ ranges from $10^3$ to $10^9$;
same-$N$ configurations with differing $T/S$ ratios isolate each axis's effect.

\paragraph{Results.}
Table~\ref{tab:ex2_grid} shows the execution times.

\paragraph{Analysis.}
At $N = 10^3$, 0.009 seconds; at $N = 10^8$, 587--643 seconds:
for a $10^5$-fold increase in $N$,
execution time increases only by a factor of approximately $6.5 \times 10^4$.
$N = 10^9$ did not complete within 1{,}000 seconds (timeout).

Even at the same $N$, the $T/S$ ratio causes differences:
at $N = 10^7$, $T{=}100{,}000,\, S{=}100$ takes 53.1\,s
while $T{=}1{,}000,\, S{=}10{,}000$ takes 64.9\,s (1.22$\times$).
Larger $T$ enables more effective gap-skipping via the $t$-axis stride
(Lemma~\ref{lem:stride}),
while larger $S$ increases the number of spatial start points.

Overall, execution time scales approximately linearly with $N$,
and $10^8$ transactions can be processed within 10 minutes.

\begin{table}[t]
\centering
\caption{EX-2: 2D scalability.
  All cells with Bernoulli background only (no embedded patterns),
  $N = T \times S$, $\theta = 15$, $W = (24, 4)$.
  Execution time (seconds). $N = 10^9$ did not complete within 1{,}000\,s.}
\label{tab:ex2_grid}
\small
\begin{tabular}{r|rrrr}
\toprule
$T \backslash S$ & 10 & 100 & 1{,}000 & 10{,}000 \\
\midrule
      100 & 0.009 &  0.059 &  0.655 &  7.18 \\
  1{,}000 &  0.042 &  0.470 &  5.62 &  64.9 \\
 10{,}000 &  0.403 &  4.79 &  57.5 &  643 \\
100{,}000 &  4.16 &  53.1 &  587 &     $>$1000 \\
\bottomrule
\end{tabular}
\end{table}

% ==============================================================
\subsection{EX-3: Application to Real-World Data ($d{=}3$--$4$)}
\label{sec:ex3}
% ==============================================================

\paragraph{Objective.}
We apply the proposed method to real-world data from diverse domains
and confirm that the detected patterns and dense regions
yield interpretable insights consistent with domain knowledge.

\paragraph{Datasets and results.}
Two crime report datasets
(Chicago Crime~\cite{chicagocrime}, NYPD Crime~\cite{nypdcrime}),
citizen complaints (NYC 311~\cite{nyc311}),
and an earthquake catalog (USGS Earthquake~\cite{usgseq})
are used, spanning three domains and $d{=}3$--$4$ dimensions.
Time is discretized into weekly bins and spatial coordinates
into equal-width bins, converting each cell's events into a transaction.
Dataset details and quantitative results
are given in Appendix~\ref{sec:dataset_details}.

\paragraph{Findings.}
In Chicago Crime, ASSAULT + BATTERY forms chronic dense windows
in south and west side districts,
and three-item patterns (e.g., THEFT + CRIMINAL DAMAGE + BATTERY)
pinpoint compound hotspots missed by single-category analysis.
In NYPD Crime, compound co-occurrences of felonies and misdemeanors
are captured at the borough and weekly granularity.
In NYC 311, combinations of noise complaints and illegal parking
repeatedly appear in residential neighborhoods.
In USGS Earthquake, 4-dimensional analysis including the depth axis
distinguishes dense patterns of shallow crustal earthquakes
from intermediate-depth intraslab events.
In all cases, the proposed method outputs
\emph{which} itemsets are dense and \emph{when and where}
as concrete spatio-temporal windows,
providing interpretable results for domain experts.
Detailed case studies are given in Appendix~\ref{sec:real_data_qualitative}.
